%!TEX program = xelatex
% ============================================================
%  main_ja.tex
%  Unity 環境を用いた MediaPipe 姿勢推定のバイアス評価
%
%  Overleaf コンパイル設定: XeLaTeX
%  初期設定は paper/paper.md（xeCJK + Noto Serif CJK JP）に準拠
%
%  Overleaf アップロード構成（source/ フォルダをそのままアップロード）:
%    main_ja.tex                      <- これをコンパイル
%    figs/                            <- main.tex と共通
% ============================================================

\documentclass[twocolumn]{article}

% ---- 日本語設定（paper.md 準拠）--------------------------
\usepackage{xeCJK}
\setCJKmainfont{Noto Serif CJK JP}

% ---- 数式・記号 -------------------------------------------
\usepackage{amsmath}
\usepackage{gensymb}

% ---- レイアウト（paper.md 準拠）--------------------------
\usepackage{geometry}
\geometry{top=24mm, bottom=24mm, left=19mm, right=19mm}
\setlength{\parskip}{0.15em}
\setlength{\floatsep}{6pt}
\setlength{\textfloatsep}{8pt}

% ---- 図・表 -----------------------------------------------
\usepackage{graphicx}
\usepackage{subcaption}
\usepackage{booktabs}
\usepackage{multirow}
\usepackage{array}
\usepackage{float}

% ---- その他（paper.md 準拠）------------------------------
\usepackage{url}
\usepackage{titlesec}
\usepackage[hidelinks]{hyperref}
\usepackage{siunitx}

\titlespacing*{\section}      {0pt}{2ex plus 0.5ex minus 0.2ex}{1.2ex plus 0.2ex}
\titlespacing*{\subsection}   {0pt}{1.8ex plus 0.5ex minus 0.2ex}{0.9ex plus 0.2ex}
\titlespacing*{\subsubsection}{0pt}{1.5ex plus 0.3ex minus 0.2ex}{0.6ex plus 0.2ex}

\sloppy

% ---- タイトルブロック -------------------------------------
\title{Unity 環境を用いた MediaPipe 姿勢推定のバイアス評価}

\author{%
  非会員\quad 平良\ 文磨$^{*a)}$ \qquad
  正員\quad   加藤\ 司$^{*b)}$ \qquad
  非会員\quad 安富祖\ 仁$^{*c)}$\footnotemark[1]
}
\date{}

% ============================================================
\begin{document}
\maketitle

\begingroup
\renewcommand{\thefootnote}{}
\footnotetext[1]{$^{*a)}$ 沖縄県立開邦高等学校，$^{*b)}$ 琉球大学大学院教育学研究科，$^{*c)}$ 株式会社lollol}
\endgroup

% ============================================================
% アブストラクト
% ============================================================
\section*{アブストラクト}

本研究は，Unityシミュレーション環境で収集したグラウンドトゥルース（GT）を
用いて，MediaPipeによる姿勢推定の\textbf{系統的バイアス}を定量的に解明することを
目的とする．従来の評価研究は正面・限定視点での精度指標（単一数値）の報告に留まっていたが，
本研究は実用場面で生じる\textbf{予測可能・構造的な誤差パターン}――
モデルアーキテクチャ，学習データ分布，座標系の慣例に起因するバイアス――を
体系的に特定・定量化する．

505台の多視点カメラ（高さ4層）から107フレームを収集し，
12関節 $\times$ 259,356件の観測データを得た．
評価指標として関節角度誤差（MAE）と方向角誤差（$\Delta\theta$，$\Delta\psi$）を用い，
以下の4種のバイアスを同定した：

\begin{itemize}\itemsep0pt
  \item \textbf{座標系バイアス}：UnityとMediaPipeのY軸方向の不整合により，
        肘の$|\Delta\theta|$は修正前$\approx121\degree{}$から
        修正後$58\degree{}$に，肩は$\approx170\degree{}$から$\approx10\degree{}$に低減した．
  \item \textbf{骨盤回転アーティファクト}：左右股関節の奥行き誤差が
        強い負の相関（$r=-0.84$）を示し，
        骨盤剛性制約の欠如による疑似的な骨盤回転が生じていた．
  \item \textbf{上肢連動バイアス}：肩・肘・手首の奥行き方向誤差が
        正の相関（$r=0.72$--$0.77$）を示し，
        運動学的連鎖に沿った誤差伝播が確認された．
  \item \textbf{視点バイアス}：俯瞰カメラ（Y$=$2.0\,m）では
        全関節のMAEが最大15\%増加し，
        BlazePoseの学習データにおける正面視点偏重を裏付けた．
\end{itemize}

\vspace{-0.8em}

% ============================================================
% 1. はじめに
% ============================================================
\section{はじめに}

姿勢推定はスポーツ分析，臨床リハビリテーション，人間工学，
ヒューマンコンピュータインタラクションの中核技術として定着している
\cite{mediapipe,blazepose}．
Googleが開発したMediaPipeは，GPUを必要とせず
コモディティハードウェア上でリアルタイム処理を実現できる点で
実務上の採用が急速に拡大している\cite{mediapipe}．

MediaPipeの3D精度に関する系統的・多視点評価は依然として不足している．
Bazarevsky~ら\cite{blazepose}はBlazePoseのベンチマーク性能を報告しているが，
その評価は視点多様性が限定されたキュレーションベンチマークに限定されている．
実際の展開環境ではカメラ配置を自由に選択できないことも多く，
視点によって精度がどう変化するかを理解することは不可欠である．

BlazePoseおよび関連推定器の独立評価研究（第2節で詳述）は
標準ベンチマーク・正面近傍視点での精度評価に集中しており，
多様なカメラ高さと水平位置にわたる系統的な多視点評価は
実施されていない．
シミュレーションベースのGT収集――SURREAL\cite{surreal}が先駆けとなり，
その後Unityベースのジェネレータ\cite{peoplesanspeople}に発展――は，
追加の計測機器なしに任意の視点から完全に制御されたカメラ配置で
大規模データを取得できるスケーラブルな手法を提供する．

単眼3D姿勢推定の根本的課題は奥行き曖昧性である：
1枚の2D画像投影から3D関節位置を一意に決定することはできず，
推定精度は学習データの視点分布に強く依存する\cite{mehta2017}．
COCO\cite{coco}をはじめとする公開データセットは
正面・立位・歩行姿勢が支配的であり，
これが学習モデルに正面視点バイアスを埋め込み，
側方・俯瞰カメラへの汎化性能を低下させる．
本研究はこの仮説を505カメラ位置で直接検証し，
さらにUnityとMediaPipeの間に存在する系統的なY軸座標不整合を同定する．
この不整合を未修正のまま放置すると方向角誤差が最大$\approx120\degree{}$膨張する．

本論文の貢献は以下のとおりである：
\begin{enumerate}\itemsep0pt
  \item \textbf{MediaPipeバイアス分類}：座標不整合，骨盤アーティファクト，
        上肢連動，視点劣化という4種の系統的バイアスを同定・定量化する．
  \item \textbf{座標系バイアス修正}：UnityとMediaPipeのY軸方向不整合を同定し，
        関節ごとの影響量（$|\Delta\theta|$の最大削減161\degree{}）を定量化し，
        修正式 $y \leftarrow -y$ を提供する．
  \item \textbf{方向角指標の導入}：$\Delta\theta$と$\Delta\psi$を
        関節角度MAEの補完指標として導入する．
        これらの指標は骨盤アーティファクトなど
        MAEでは不可視のバイアスパターンを顕在化させる．
  \item \textbf{大規模多視点データセット}：505カメラ位置 $\times$
        107フレームによる259,356件の観測データを構築し，
        標準ベンチマークにない側方・俯瞰視点を含む
        初の系統的視点バイアス分析を実現する．
\end{enumerate}

% ============================================================
% 2. 関連研究
% ============================================================
\section{関連研究}

\subsection{MediaPipeとBlazePoseの精度評価}

BlazePoseはBazarevsky~ら\cite{blazepose}によって提案され，
COCOバリデーションセットでPCK@0.2 = 72.0\,\%，
独自のフィットネス動画ベンチマークでState-of-the-artを達成したと報告されている．
ただしその評価は，フィットネス映像に典型的な正面に近い視点に限定されており，
多様なカメラ視点での性能は評価されていない．

BlazePoseの元評価\cite{blazepose}は多様なカメラ視点にわたる
系統的な精度データを提供していない．
本研究はこのギャップを埋め，
高さ4層（Y$=$0.5--2.0\,m）を網羅する505カメラの3Dグリッドで
系統的な視点別精度分析を初めて実施する．

\subsection{シミュレーションベースのGT収集}

実世界で大規模かつ正確な3D GTを得ることは困難である：
手動3Dアノテーションは労働集約的であり，
光学式モーションキャプチャシステムは高コストかつ実験室限定である．
合成レンダリングはその代替として標準的アプローチとなっている．
Varol~ら\cite{surreal}はSURREALを提案し，
SMPLボディアバターの完全な2D/3D姿勢・深度・セグメンテーションラベルを
600万フレーム以上生成した．
Patel~ら\cite{agora}はAGORAで，
高品質テクスチャを持つボディスキャンを自然な屋外環境に合成し，
現実的な合成データが3DPWベンチマークの性能を向上させることを示した．

ゲームエンジンベースのシステムはカメラ完全制御という追加利点を持つ．
Ebadi~ら\cite{peoplesanspeople}はUnityベースの合成データジェネレータ
PeopleSansPeopleを公開した．
本研究も同様にUnityシミュレーション環境を活用し，
505台の系統的に配置されたカメラから正確なフレーム単位3D関節座標を記録する
——物理的なモーションキャプチャ設備では経済的に実現困難な視点一斉取得である．

\subsection{単眼3D姿勢推定と奥行き曖昧性}

単眼RGBカメラから3D姿勢を復元することは本質的に不良設定問題である：
多くの3D配置が透視投影で同一の2D画像に写像される．
Mehta~ら\cite{mehta2017}はマルチデータセット転移学習によって野外シーンへの
汎化を改善し，標準的なポーズベンチマークより視点多様性の高い
MPI-INF-3DHPベンチマークを導入したが，
側方・俯瞰視点は学習コーパスでなお過少代表されたままである．

COCO\cite{coco}をはじめとする公開データセットは
正面に近い立位・歩行姿勢が支配的であり，
学習モデルに強い正面視点バイアスを埋め込む．
このバイアスにより，側方・俯瞰配置のカメラでは精度が低下する
——本研究の多層評価がこれを定量的に確認する（表~\ref{tab:mae_by_layer}）．
また，MediaPipeの骨格モデルは骨盤をランドマーク23（左股関節）と
24（右股関節）を結ぶ単純な線分として表現しており，
剛体多セグメント構造を持たない．
このため2D画像上のわずかな水平画素差が3D復元時に
大きな骨盤回転として誤解釈され，
第4節で報告する左右対称な股関節奥行き誤差（$r=-0.840$）を生じさせる．

\subsection{骨格拘束モデル}

奥行き曖昧性を緩和するために解剖学的・時間的拘束を課す
アーキテクチャが提案されている．
Mehta~ら\cite{vnect}は骨長一貫性とフレーム間時間安定性を課す
運動学的骨格フィッティング段階を持つリアルタイム単眼3D姿勢推定を実現した．
PoseFormer\cite{poseformer}は畳み込み層を持たない純粋なトランスフォーマーベースの
ビデオ姿勢推定器であり，Human3.6Mで44.3\,mm MPJPEを達成した．
MotionBERT\cite{motionbert}はノイズのある2D観測から3Dモーションを復元する
デュアルストリーム時空間トランスフォーマーを事前学習し，
Human3.6Mで39.2\,mm MPJPEを達成した．

これらの進歩にもかかわらず，骨盤剛性を陽的な硬拘束として課す手法は稀である．
本研究で観測された股関節の反相関（$r=-0.840$）は，
歩行動作における左右股関節深度を剛体骨盤モデルの下で
相互整合させる拘束が奥行き誤差を大幅に低減できる可能性を示唆しており，
将来のモデル設計に向けた推奨方向として提示する．

% ============================================================
% 3. 実験方法
% ============================================================
\section{実験方法}

\subsection{実験環境}

Unity（バージョン6000.0.60f1）を用いて，
Y-Botヒューマノイドアバターの歩行シミュレーション環境を構築した．
アバターは座標$(0,\,0,\,-3)$から$(0,\,0,\,3)$まで
Z軸正方向に直進歩行した．

\begin{figure}[!t]
\centering
\includegraphics[width=0.30\textwidth]{figs/fig01_camera_layout.png}
\caption{カメラ配置図（高さ4層 Y$=$0.5, 1.0, 1.5, 2.0\,mにわたる505台）．
  中央の矢印はアバターの歩行方向（+Z軸方向）を示す．}
\label{fig:camera_layout}
\end{figure}

カメラは高さ4層（Y$=$0.5, 1.0, 1.5, 2.0\,m）と
水平位置（X, Z $\in [-6, 6]$\,m）を網羅する3Dグリッドに配置した．
各層では $13\times13$ グリッドから中央の $5\times5$ ブロックを除外し，
理論上最大 $4(13^2-5^2)=576$ 台のうち505台でデータを収集した．
各位置から解像度 $1280\times720$ピクセルで107フレームを撮影し，
総計259,356件の関節観測値（可視性フィルタリング後）を得た．

使用したアバターはY-Botメッシュ（Unity標準ヒューマノイドリグ）に
組み込みループ歩行アニメーションクリップを適用したものである．
モーションはZ軸正方向に約6\,mのパスを歩行し，
1サイクル分の自然な関節角度変化を含む．

\begin{figure}[!t]
  \centering
  \includegraphics[width=0.44\textwidth]{figs/fig02a_unity_env.png}
  \vspace{0.3em}
  \includegraphics[width=0.44\textwidth]{figs/fig02b_frame_sample.jpg}
  \caption{Unity実験環境（上）と撮影フレーム例（下，Y$=$0.5\,m, X$=$-1.0\,m, Z$=$-3.0\,m）．}
  \label{fig:unity_env}
\end{figure}

\subsection{評価指標}

\begin{figure}[!t]
\centering
\includegraphics[width=0.48\linewidth]{figs/fig03_concept_metrics.png}
\caption{2種の評価指標の概念図．左：関節角度（MAE）——3点による屈曲角
  （$0$--$180\degree{}$）．右：方向角（$\Delta\theta$，$\Delta\psi$）——
  股関節中心を基準とした角度オフセット（$-180\degree{}$から$+180\degree{}$）．}
\label{fig:concept_metrics}
\end{figure}

\paragraph{関節角度MAE}
3点で定義される屈曲角をGTとMediaPipeそれぞれで計算し，
絶対差の平均（MAE）を求める．これは関節可動域の推定精度を表す．

\paragraph{方向角誤差（$\Delta\theta$，$\Delta\psi$）}
股関節中心を原点として各関節の3D方向ベクトルを球面角（仰角$\theta$，
方位角$\psi$）で表現し，GTとMediaPipe間のずれを求める．
MAEが捉えられないモデルの系統的な方向バイアスを定量化する．

\subsection{座標系統一}

UnityはY軸上向き，MediaPipeはY軸下向きの座標系を採用する．
このY軸方向の不整合により，未修正のMediaPipe座標から計算した$\theta$の
符号が全関節で反転する．修正式は：
\[
  y_{\text{unified}} = -y_{\text{MediaPipe}}
\]
この変換を\texttt{06\_direction\_detection/scripts/coordinate\_transform.py}
に実装した．以下で報告する全ての方向角結果はこの座標統一後の値を用いる．

% ============================================================
% 4. 実験結果
% ============================================================
\section{実験結果}

\subsection{関節角度MAE}

表~\ref{tab:joint_angle_error}に8主要関節の
MAE・中央値・最大値をまとめる．
肘と膝が最も誤差が小さく（それぞれ$\approx18\degree{}$，$\approx17$--$19\degree{}$），
肩が最大のMAEを示す（$\approx39$--$40\degree{}$）——
これはカメラに対して近位にある肩関節の奥行き曖昧性が大きいためと考えられる．

\begin{table}[!t]
\centering
\caption{関節角度誤差統計（関節角度MAE，単位：度）．
  出典：\texttt{04\_mae\_heatmap/coordinate\_angle\_mae.csv}}
\label{tab:joint_angle_error}
\begin{tabular}{lccc}
\toprule
関節          & MAE    & 中央値 & 最大値  \\
              & (deg)  & (deg)  & (deg)  \\
\midrule
左肩  & 39.8   & 39.2   & 70.1   \\
右肩  & 39.0   & 40.0   & 58.2   \\
左肘  & 18.6   & 15.9   & 56.9   \\
右肘  & 18.2   & 13.6   & 46.2   \\
左股関節  & 30.2   & 29.4   & 66.0   \\
右股関節  & 33.3   & 32.2   & 81.7   \\
左膝  & 17.0   & 16.6   & 43.4   \\
右膝  & 19.5   & 19.4   & 62.3   \\
\bottomrule
\end{tabular}
\end{table}

カメラ高さ層別のMAEは表~\ref{tab:mae_by_layer}と
図~\ref{fig:mae_layer_bar}に示す．
図~\ref{fig:mae_layer_compare}は左肘の3層ヒートマップを示す．
Y$=$2.0\,m（俯瞰）で誤差が顕著に増加しており，
BlazePoseの学習データに埋め込まれた正面視点バイアスと整合する．

\subsection{方向角誤差（$\Delta\theta$，$\Delta\psi$）}

座標系統一後の全12関節の方向角誤差統計を
表~\ref{tab:direction_angle_error}に示す．

\begin{table}[!t]
\centering
\caption{座標統一後の方向角誤差統計（単位：度，絶対値平均）．
  出典：\texttt{11\_direction\_ditection/output/processed\_data/joint\_summary.csv}．
  $^\ddagger$股関節の値は\texttt{detailed\_results.csv}から集計
  （$n=21{,}613$件/関節，\texttt{joint\_summary.csv}から除外されているため）．}
\label{tab:direction_angle_error}
\small
\begin{tabular}{lcccc}
\toprule
関節 & \multicolumn{2}{c}{$|\Delta\theta|$} & \multicolumn{2}{c}{$|\Delta\psi|$} \\
\cmidrule(lr){2-3}\cmidrule(lr){4-5}
     & 平均 & SD & 平均 & SD \\
\midrule
左肩   & 10.4 & 11.7 & 94.1 & 86.4 \\
右肩   &  9.3 & 10.1 & 92.7 & 92.0 \\
左肘   & 58.3 & 24.1 & 87.3 & 90.6 \\
右肘   & 59.3 & 22.7 & 88.6 & 94.5 \\
左手首 & 84.3 & 103.7 & 88.6 & 91.9 \\
右手首 & 91.5 & 111.9 & 91.6 & 99.2 \\
左股関節$^\ddagger$ & 83.5 & 72.5 & 88.9 & 26.4 \\
右股関節$^\ddagger$ & 96.5 & 72.5 & 91.1 & 26.4 \\
左膝   & 17.4 & 18.8 & 84.7 & 97.9 \\
右膝   & 18.5 & 19.1 & 87.2 & 97.5 \\
左足首 & 11.3 & 14.0 & 90.6 & 105.8 \\
右足首 & 11.9 & 14.5 & 100.4 & 103.1 \\
\bottomrule
\end{tabular}
\end{table}

図~\ref{fig:elbow_heatmaps}に肘$\Delta\theta$のカメラ位置別分布を示す．

\begin{figure}[!t]
  \centering
  \begin{subfigure}[b]{0.48\linewidth}
    \includegraphics[width=\linewidth]{figs/fig04a_elbow_L_theta_y05.jpg}
    \caption{左肘 $\Delta\theta$（Y$=$0.5）}
  \end{subfigure}\hfill
  \begin{subfigure}[b]{0.48\linewidth}
    \includegraphics[width=\linewidth]{figs/fig04b_elbow_R_theta_y05.jpg}
    \caption{右肘 $\Delta\theta$（Y$=$0.5）}
  \end{subfigure}
  \caption{肘方向角誤差$\Delta\theta$のヒートマップ（Y$=$0.5）．
    左右肘が逆符号パターンを示しており，
    強い負の相関（$r=-0.862$）と整合する．}
  \label{fig:elbow_heatmaps}
\end{figure}

図~\ref{fig:hip_psi_heatmap}に股関節$\Delta\psi$のヒートマップを示す．
左右股関節はほぼ同一の絶対値パターンを示し，
強い負の相関（$r=-0.840$）による骨盤回転アーティファクトを視覚化する．

\begin{figure}[!t]
  \centering
  \begin{subfigure}[b]{0.48\linewidth}
    \includegraphics[width=\linewidth]{figs/fig05a_hip_L_psi_y05.jpg}
    \caption{左股関節 $\Delta\psi$（Y$=$0.5）}
  \end{subfigure}\hfill
  \begin{subfigure}[b]{0.48\linewidth}
    \includegraphics[width=\linewidth]{figs/fig05b_hip_R_psi_y05.jpg}
    \caption{右股関節 $\Delta\psi$（Y$=$0.5）}
  \end{subfigure}
  \caption{股関節方向角誤差$\Delta\psi$のヒートマップ（Y$=$0.5）．
    $\Delta\psi_{\text{L}} \approx -\Delta\psi_{\text{R}}$（$r=-0.840$）のため，
    左右がほぼ逆転した対称パターンを示す．}
  \label{fig:hip_psi_heatmap}
\end{figure}

\begin{figure*}[t]
  \centering
  \begin{subfigure}[b]{0.31\textwidth}
    \includegraphics[width=\linewidth]{figs/fig07a_mae_elbow_L_y05.png}
    \caption{Y$=$0.5}
  \end{subfigure}\hfill
  \begin{subfigure}[b]{0.31\textwidth}
    \includegraphics[width=\linewidth]{figs/fig07b_mae_elbow_L_y15.png}
    \caption{Y$=$1.5}
  \end{subfigure}\hfill
  \begin{subfigure}[b]{0.31\textwidth}
    \includegraphics[width=\linewidth]{figs/fig07c_mae_elbow_L_y20.png}
    \caption{Y$=$2.0}
  \end{subfigure}
  \caption{左肘の関節角度MAEヒートマップ（カメラ高さ3層）．
    Y$=$2.0（俯瞰）で誤差が増加し，学習データの視点バイアスと整合する．
    出典：\texttt{4\_MAE\_HEATMAP/}．}
  \label{fig:mae_layer_compare}
\end{figure*}

\begin{figure}[!t]
  \centering
  \includegraphics[width=\linewidth]{figs/fig08_mae_layer_bar.png}
  \caption{カメラ高さ層別（Y$=$0.5, 1.0, 1.5, 2.0\,m）の関節角度MAE．
    全関節でY$=$2.0で誤差が増加し，右股関節で最大（Y$=$1.0比$+4.7\degree{}$，$+15\%$）．
    出典：\texttt{4\_MAE\_HEATMAP/}．}
  \label{fig:mae_layer_bar}
\end{figure}

\begin{table}[!t]
\centering
\caption{カメラ高さ層別の関節角度MAE（単位：度）．
  出典：\texttt{4\_MAE\_HEATMAP/Y=0.5,1.5/} および
  \texttt{4\_MAE\_HEATMAP/Y=1.0.2.0/coordinate\_angle\_mae.csv}．}
\label{tab:mae_by_layer}
\small
\begin{tabular}{lcccc}
\toprule
関節 & Y$=$0.5 & Y$=$1.0 & Y$=$1.5 & Y$=$2.0 \\
\midrule
左肩  & 39.8 & 38.9 & 39.6 & 41.0 \\
右肩  & 39.3 & 38.1 & 38.6 & 40.2 \\
左肘  & 18.1 & 17.7 & 18.3 & 20.2 \\
右肘  & 17.5 & 16.7 & 18.2 & 20.2 \\
左股関節 & 29.6 & 29.0 & 29.9 & 32.5 \\
右股関節 & 32.0 & 31.7 & 32.8 & 36.7 \\
左膝  & 16.1 & 16.6 & 17.7 & 17.8 \\
右膝  & 18.9 & 18.7 & 19.6 & 20.6 \\
\bottomrule
\end{tabular}
\end{table}

\subsection{関節間誤差相関}

表~\ref{tab:correlation_theta}と表~\ref{tab:correlation_psi}に
$\Delta\theta$および$\Delta\psi$の高相関ペア（$|r|>0.7$）を示す．
全ペアで$p<0.001$（観測数$n>21{,}000$フレーム--カメラペア）．

\begin{table}[!t]
\centering
\caption{XY平面（$\Delta\theta$）における高相関関節ペア（$|r|>0.7$）．
  全ペア$p<0.001$（$n>21{,}000$件）．}
\label{tab:correlation_theta}
\small
\begin{tabular}{llr}
\toprule
関節1        & 関節2        & $r$ \\
\midrule
左肘    & 右肘   & $-0.862$ \\
左肘    & 左肩   &  $0.788$ \\
右足首  & 右膝   &  $0.767$ \\
左足首  & 左膝   &  $0.766$ \\
右肘    & 右肩   &  $0.710$ \\
\bottomrule
\end{tabular}
\end{table}

\begin{table}[!t]
\centering
\caption{XZ平面（$\Delta\psi$）における高相関関節ペア（$|r|>0.7$）．
  全ペア$p<0.001$（$n>21{,}000$件）．}
\label{tab:correlation_psi}
\small
\begin{tabular}{llr}
\toprule
関節1        & 関節2         & $r$ \\
\midrule
左股関節  & 右股関節  & $-0.840$ \\
右肘      & 右肩      &  $0.770$ \\
右肘      & 右手首    &  $0.769$ \\
左肘      & 左肩      &  $0.768$ \\
右肩      & 右手首    &  $0.726$ \\
左肘      & 左手首    &  $0.722$ \\
右肘      & 右股関節  &  $0.721$ \\
左股関節  & 右肘      & $-0.705$ \\
\bottomrule
\end{tabular}
\end{table}

\begin{figure*}[t]
  \centering
  \begin{subfigure}[b]{0.49\textwidth}
    \centering
    \includegraphics[width=\linewidth]{figs/fig06a_corr_theta.png}
    \caption{$\Delta\theta$（XY平面）}
  \end{subfigure}\hfill
  \begin{subfigure}[b]{0.49\textwidth}
    \centering
    \includegraphics[width=\linewidth]{figs/fig06b_corr_psi.png}
    \caption{$\Delta\psi$（XZ平面）}
  \end{subfigure}
  \caption{関節間方向角誤差の相関行列．
    (a) $\Delta\theta$は表~\ref{tab:correlation_theta}，
    (b) $\Delta\psi$は表~\ref{tab:correlation_psi}に対応する．
    出典：\texttt{06\_direction\_detection/output/correlation\_analysis/}．}
  \label{fig:correlation_heatmaps}
\end{figure*}

\subsection{3D位置誤差（MPJPE）}

補足的な位置指標として，推定3Dランドマーク位置とGT位置の
ユークリッド距離（MPJPE：Mean Per-Joint Position Error）を算出した．
MediaPipe world座標系（身長1.7\,mの人物で1単位$\approx$1\,m）における
全12関節の全体平均は0.363\,mである．
股関節は座標系の原点であるため誤差が小さく（0.044\,m），
遠位関節では誤差が大きい：足首（0.71\,m），
手首（0.36--0.40\,m），肘（0.31--0.32\,m）．
足首の大きな誤差は単眼復元における奥行き軸不確かさを主に反映しており，
股関節原点からの誤差が運動学的連鎖に沿って蓄積する．

\subsection{カメラ視点依存性}

フレームごとに全関節平均$|\Delta\theta|$が最小・最大となるカメラを同定した
（\texttt{frame\_camera\_summary.csv}，107フレーム$\times$201カメラ/フレーム）．

\emph{最良}カメラは全フレーム平均で$|\Delta\theta|=8.8\degree{}$
（最小6.1\degree{}，最大15.8\degree{}）を達成し，
最適視点が存在することを示す．
一方，\emph{最悪}カメラは平均80.1\degree{}
（最小55.6\degree{}，最大148.1\degree{}）となり，
最悪視点ではランダム方向推定に匹敵する誤差が生じることを示す．
同一ボーズに対する最良・最悪カメラの9倍差は，
実用展開においてカメラ配置が如何に重要かを際立たせる．

% ============================================================
% 5. 考察
% ============================================================
\section{考察}

\subsection{座標系統一効果}

UnityとMediaPipeのY軸方向不整合は全関節の$\theta$に系統的な符号誤差を生じさせる．
Y軸が反転している場合$\theta_{\text{MP}} \approx -\theta_{\text{GT}}$となるため，
未修正の方向角誤差は$\Delta\theta \approx -2\theta_{\text{GT}}$となる；
$\theta_{\text{GT}} \approx 60\degree{}$の関節では$|\Delta\theta| \approx 120\degree{}$に達する．
$y \leftarrow -y$を適用後，\emph{肩}の平均$|\Delta\theta|$は
$\approx120\degree{}$から$\approx10\degree{}$に，
\emph{肘}は$\approx121\degree{}$から$\approx58\degree{}$に低減した\cite{mediapipe_coord_issue}．

表~\ref{tab:yflip_effect}に関節ごとの改善量を示す．
肩・足首は最も恩恵を受け（$\Delta|\Delta\theta|>155\degree{}$），
手首は最小の改善に留まる（$<10\degree{}$）．
これは手首の推定分散が大きく，系統的な符号誤差を雑音が覆い隠すためである．

\begin{table}[!t]
\centering
\caption{Y軸座標統一の効果：$y \leftarrow -y$適用前後の平均$|\Delta\theta|$（単位：度）．
  出典：\texttt{7\_direction\_ditection}（修正前），
  \texttt{11\_direction\_ditection}（修正後）．}
\label{tab:yflip_effect}
\small
\begin{tabular}{lrrr}
\toprule
関節     & 修正前 & 修正後 & 削減量 \\
\midrule
左肩  & 169.5 & 10.4 & 159.1 \\
右肩  & 171.2 &  9.3 & 161.9 \\
左肘  & 121.6 & 58.3 &  63.3 \\
右肘  & 121.2 & 59.3 &  61.9 \\
左手首 &  94.3 & 84.3 &  10.0 \\
右手首 &  93.1 & 91.5 &   1.6 \\
左膝  & 161.7 & 17.4 & 144.3 \\
右膝  & 162.2 & 18.5 & 143.7 \\
左足首 & 168.7 & 11.3 & 157.4 \\
右足首 & 167.9 & 11.9 & 156.0 \\
\bottomrule
\end{tabular}
\end{table}

統一後も残存する$\approx58\degree{}$の肘誤差は，
座標系の不整合ではなく，学習データの視点不均衡と
「腕を体側に下げた」姿勢バイアスによる真のバイアスを反映する．

\subsection{奥行き推定における系統的バイアス}

MediaPipeのBlazePoseは各ランドマークをヒートマップと回帰で推定し，
股関節中心を基準とした3D world座標を出力する\cite{blazepose}．
このアーキテクチャは骨長一貫性や骨盤剛性などの
解剖学的拘束を陽的に組み込んでいない．
座標統一後も肘$|\Delta\theta|\approx58\degree{}$の残差が持続し，
方向バイアスが根強く残ることを示す．

COCO\cite{coco}などの学習データセットは
日常的・正面近傍・立位/歩行姿勢が支配的であり，
肘方向推定を下垂姿勢（体側）にバイアスさせる可能性が高い．
左右肘誤差の逆符号パターン（$r=-0.862$）は
学習時の左右フリップデータ拡張を強く示唆し，
肘$\Delta\theta$のSDが$\approx23$--$25\degree{}$と小さいことは，
このバイアスが比較的予測可能で後処理補正が可能であることを示す．

\subsection{左右対称な股関節奥行き誤差}

股関節では$|\Delta\psi|\approx90\degree{}$かつ強い負の相関（$r=-0.840$）が示すように，
奥行き方向誤差が反相関を示す：一方の股関節が前方に誤推定されると
他方が後方に誤推定され，疑似骨盤回転アーティファクトが生じる．
左右股関節のヒートマップが同一の絶対値パターンを示すことは
（図~\ref{fig:hip_psi_heatmap}），
各観測点での$\Delta\psi_{\text{L}} \approx -\Delta\psi_{\text{R}}$と整合する．

MediaPipeの骨格モデルでは骨盤がランドマーク23（左股関節）と
24（右股関節）を結ぶ線分として表現されており，
仙骨・腸骨のような多セグメント剛体構造を持たない．
したがって3D復元時の2つの股関節の相対深度に剛性拘束が課されない．
その結果，2D画像上のわずかな水平画素差が
一方の股関節を近く，他方を遠く配置することを意味する疑似骨盤回転として
解釈され，左右対称誤差パターンが生じる．
骨盤剛性を硬拘束または軟拘束として課すモデル改善が必要である．

\subsection{上肢連動バイアス}

$\Delta\psi$における肩・肘・手首の高い正の相関（$r=0.72$--$0.77$）は，
誤差が運動学的連鎖に沿って階層的に伝播しているか，
あるいは3D復元時に上肢が暗黙的に剛体として扱われていることを示唆する．
グラフニューラルネットワークや陽的運動学モデルによる
骨長拘束の組み込みがこの伝播を低減できる可能性がある．

\subsection{カメラ視点依存性}

フレームごとの最良カメラを選択すると$|\Delta\theta|$は
6.1\degree{}--15.8\degree{}（平均8.8\degree{}）に低減するが，
最悪カメラでは55.6\degree{}--148.1\degree{}（平均80.1\degree{}）となる．
この9倍の差は推定品質が視点に強く依存することを確認する．
BlazePoseの学習データ（主にCOCO\cite{coco}およびフィットネス動画\cite{blazepose}）は
正面・正面近傍視点が支配的であり，
極端な角度（側方・俯瞰）のカメラは過少代表されていることで汎化性能が低下する．
これは表~\ref{tab:mae_by_layer}でY$=$2.0（俯瞰カメラ）で
MAEが上昇する観測と整合する．

\subsection{実用的示唆}

MediaPipeは2D関節位置推定に有用であるが，
3D奥行き推定は単眼曖昧性により根本的に限界がある．
画像平面内の相対関節角度で十分なアプリケーション
（水泳のストロークカウント分析，粗大運動リハビリスクリーニングなど）では
MediaPipeで十分対応できる．
3D身体方向を必要とするアプリケーション
（歩容分析，体幹回転の人間工学リスク評価など）では，
ここで同定したバイアス――特に股関節$\Delta\psi$と肘$\Delta\theta$――を
補償する必要がある．

\subsection{本研究の限界}

本研究には以下の限界がある：
（1）アバターはY-Botメッシュの標準設定であり，
実被験者に存在する皮膚変形・衣服の効果は含まれない；
（2）単一の歩行動作のみを評価したため，
他の動作パターンでは異なる誤差パターンが生じる可能性がある；
（3）照明は均一なシミュレーション環境であり，
実世界の照明変動は評価していない；
（4）MediaPipeのみを評価しており，
他の推定器（YOLO-Pose，OpenPoseなど）との比較は今後の課題である．

% ============================================================
% 6. 結論
% ============================================================
\section{結論}

本研究では，Unityシミュレーション環境を用いた
MediaPipe姿勢推定の大規模多視点バイアス評価を実施した．
主な知見を以下に示す：

\begin{itemize}\itemsep0pt
  \item 関節角度MAEは$\approx17\degree{}$（膝）から$\approx40\degree{}$（肩）の範囲にあり，
        肘と膝が最も誤差が小さい．
        Y$=$2.0（俯瞰）カメラではMAEが増加し，学習データ視点バイアスと整合する．
  \item 座標系統一（Y軸反転）は必須である：
        肩$|\Delta\theta|$を$\approx120\degree{}$から$\approx10\degree{}$に，
        肘$|\Delta\theta|$を$\approx121\degree{}$から$\approx58\degree{}$に削減する．
  \item 肘の方向誤差（$|\Delta\theta|\approx58\degree{}$）は
        学習データバイアスを反映し，系統的（SD $\approx23$--$25\degree{}$）であるため，
        後処理補正の余地がある．
  \item 股関節奥行き誤差（$|\Delta\psi|\approx90\degree{}$，$r=-0.84$）は
        骨格モデルの骨盤剛性拘束欠如を反映する．
  \item 上肢誤差は強く連動し（$r=0.72$--$0.77$），
        階層的誤差伝播を示唆する．
\end{itemize}

将来の方向性として：
（1）姿勢推定モデルへの骨盤剛性拘束の組み込み；
（2）多様な角度の多視点データを用いたファインチューニング；
（3）他推定器・動作タイプへの評価フレームワークの拡張；
（4）肘・股関節誤差に対するアプリケーション特化補正フィルタの開発
が挙げられる．

% ============================================================
\section*{謝辞}
% ============================================================

本研究は，科学技術振興機構（JST）
次世代人材育成事業（グローバルサイエンスキャンパス）の支援を受けた．

% ============================================================
% 参考文献
% ============================================================
\begin{thebibliography}{99}

\bibitem{mediapipe}
Google LLC, ``MediaPipe Pose Landmarker,''
\url{https://ai.google.dev/edge/mediapipe/solutions/vision/pose_landmarker},
2024.

\bibitem{mediapipe_coord_issue}
Google, ``Orientation of pose world landmarks,''
GitHub Issue \#3370, MediaPipe repository, 2022.
\url{https://github.com/google/mediapipe/issues/3370}

\bibitem{blazepose}
V.~Bazarevsky, I.~Grishchenko, K.~Raveendran, T.~Zhu, F.~Zhang,
and M.~Grundmann,
``BlazePose: On-device real-time body pose tracking,''
\textit{arXiv preprint arXiv:2006.10204}, 2020.

\bibitem{unity}
Unity Technologies, ``HumanBodyBones,''
\textit{Unity Scripting Reference}, 2024,
\url{https://docs.unity3d.com/ScriptReference/HumanBodyBones.html}.

\bibitem{coco}
T.-Y.~Lin et al.,
``Microsoft COCO: Common Objects in Context,''
in \textit{Proc. ECCV}, 2014, pp.~740--755.

\bibitem{openpose}
Z.~Cao, G.~Hidalgo, T.~Simon, S.-E.~Wei, and Y.~Sheikh,
``OpenPose: Realtime Multi-Person 2D Pose Estimation Using Part Affinity
Fields,''
\textit{IEEE Trans. Pattern Anal. Mach. Intell.}, vol.~43, no.~1,
pp.~172--186, 2021.

\bibitem{surreal}
G.~Varol et al.,
``Learning from synthetic humans,''
in \textit{Proc. CVPR}, 2017, pp.~109--117.

\bibitem{agora}
P.~Patel et al.,
``AGORA: Avatars in geography optimized for regression analysis,''
in \textit{Proc. CVPR}, 2021.

\bibitem{peoplesanspeople}
S.~E.~Ebadi et al.,
``PeopleSansPeople: A synthetic data generator for human-centric
computer vision,''
\textit{arXiv preprint arXiv:2112.09290}, 2021.

\bibitem{mehta2017}
D.~Mehta et al.,
``Monocular 3D human pose estimation in the wild using improved CNN
supervision,''
in \textit{Proc. Int.\ Conf.\ 3D Vision (3DV)}, 2017, pp.~506--516.

\bibitem{vnect}
D.~Mehta et al.,
``VNect: Real-time 3D human pose estimation with a single RGB camera,''
\textit{ACM Trans.\ Graph.}, vol.~36, no.~4, pp.~44:1--44:14, 2017.

\bibitem{poseformer}
C.~Zheng et al.,
``3D human pose estimation with spatial and temporal transformers,''
in \textit{Proc. ICCV}, 2021, pp.~11656--11666.

\bibitem{motionbert}
W.~Zhu et al.,
``MotionBERT: A unified perspective on learning human motion
representations,''
in \textit{Proc. ICCV}, 2023.

\end{thebibliography}

\end{document}
